Este trabajo cumplió con el objetivo general de diseñar un kit educativo para la clasificación de objetos mediante visión por computador empleando el robot PhantomX Pincher, la plataforma Raspberry Pi 5 y el sistema operativo robótico ROS2. Los objetivos específicos se alcanzaron completamente.

Se identificaron los requerimientos de hardware y software necesarios para el sistema de clasificación, seleccionando componentes que equilibran capacidad computacional, compatibilidad con ROS2 y viabilidad económica. El stack tecnológico integra herramientas estándar de la industria para robótica de servicio y visión artificial.

Se diseñó la estructura física del kit contemplando aspectos de funcionalidad, modularidad y facilidad de manufactura. El diseño incluye plataforma base, zonas de trabajo para clasificación, sistema de captura de imágenes, y componentes de soporte estructural. La documentación técnica generada permite reproducción del diseño mediante manufactura aditiva y técnicas convencionales de mecanizado.

Se propuso una arquitectura de software modular basada en ROS2 que separa las funcionalidades de percepción, planificación y control. La arquitectura permite transición fluida entre entornos de simulación y operación con hardware real, facilitando el proceso de desarrollo y pruebas.

Se desarrollaron capacidades de simulación que permiten validar el funcionamiento del sistema sin requerir hardware físico. El workspace de ROS2 se documentó completamente y se acompañó de material educativo para facilitar su adopción en contextos académicos.

Se establecieron lineamientos pedagógicos para implementación del kit en entornos educativos, considerando aspectos de progresión curricular, evaluación de aprendizajes, y estrategias para abordar dificultades comunes en la enseñanza de sistemas robóticos.

Se construyeron prototipos físicos del kit en diferentes configuraciones según disponibilidad de componentes, validando la viabilidad del diseño propuesto y su proceso de manufactura. Los prototipos demuestran la factibilidad de producir kits educativos funcionales con recursos de laboratorio universitario.

Los entregables digitales incluyen repositorios con imagen del sistema operativo preconfigurado, modelado tridimensional completo, y código fuente documentado. Esta documentación permite que otras instituciones educativas puedan replicar o adaptar el diseño según sus necesidades específicas.

El desarrollo del kit mediante diseño personalizado resulta significativamente más económico que la adquisición de soluciones comerciales equivalentes, demostrando la viabilidad de crear herramientas educativas de calidad sin comprometer funcionalidad.